\scanpage{4}
\uline{L'icosaèdre}

$E$ espace vectoriel euclidien de dim $3$\inserted{[espace vectoriel muni d'une forme quadratique non dégénérée, qu'on supposera définie positive au besoin]}, produit scalaire $x.y$ ($\|x\|=\sqrt{x^2}$).

$S$ la sphère unité $=\{x\in E\mid\|x\|^2=1\}$.

Si $u\in S$, $u^\perp$ désigne le plan orthogonal ; donc $u+u^\perp$ est le plan tangent à $S$ en $u$, tandis que l'espace vectoriel tangent à $u$ en $S$ est can. isom. à $u^\perp$

L'ensemble des directions tangentes en $u$ \deleted{sur} \inserted{à} la sphère \inserted{$S$} \unsure{s'identifie|s'identifient} à l'ensemble $u^\perp\cap S=\Sigma_u$ des vecteurs unité orthogonaux à $u$.

On suppose $E$ orienté, de sorte que \deleted{\unsure{à}} à un couple $u\in S,v\in\Sigma_u$ est associé un troisième vecteur $w=u\wedge v$, caractérisé par les conditions $w\in S\cap u^\perp\cap v^\perp$, $(u,v,w)$ orienté en sens direct.
\inserted{On peut dire que le plan $u^\perp$ est orienté : \deleted{\unsure{puisque|précisément}} la rotation d'angle $\frac{\pi}{2}$ $v\mapsto u\wedge v$ y est définie}

Le groupe spécial orthogonal $SO(E)$ opère de façon simplement transitive sur l'ensemble des directions tangentes à la sphère $S$ en ses différents points, i.e. l'ensemble des couples $(u,v)\in S$ avec $u.v=0$ i.e. $v\in\Sigma_u$. Le groupe orthogonal $O(E)\simeq\deleted{\text{\unsure{produit}}}SO(E).\deleted{\unsure{Z}}\unsure{\mu_2|\boldsymbol{\mu}_2}$ ($\mu_2={}$homothéties $\pm1$) opère de façon doublement transitive.

\deleted{\unsure{Soit les|Soit des}} \inserted{Les} points sur $S$ du grand cercle passant par $u$ dans la direction $v$, \deleted{sont} i.e. de $w^\perp\cap S$, sont de la forme
\[
x(u,v;\alpha)=u\cos\alpha+v\sin\alpha \qquad \alpha\ \text{défini mod }2\pi\mathbb Z
\]

D'autre part, les \deleted{\unsure{vecteurs|autres}} directions tangentes à $S$ en $u$, i.e. les \deleted{\unsure{k}} points de $u^\perp\cap S$, s'écrivent
\[
v'=v\cos\theta+w\sin\theta
\]

Fixons $\theta$, et calculons les carrés des côtés du triangle de sommets $u,x(u,v;\alpha),x(u,v',\alpha)$. On a

\figure{\begin{tikzpicture}[x=1cm,y=1cm,>=Stealth,scale=.86]
\coordinate (o) at (1.45,.95); \coordinate (u) at (1.95,2.65);
\coordinate (x) at (.46,3.38); \coordinate (y) at (.42,2.04);
\coordinate (v) at (.27,1.45); \coordinate (p) at (-.10,.70); \coordinate (w) at (.18,-.04);
\draw[->,thick] (o)--(u) node[right] {$u$};
\draw[->] (o)--(v) node[left] {$v$}; \draw (o)--(p) node[left] {$v'$};
\draw[->,thick] (o)--(w) node[below] {$w$};
\draw (v) to[bend right=27] (p) to[bend right=17] (w);
\draw[dashed] (o)--(x) (o)--(y);
\draw (x) node[above] {$x(u,v;\alpha)$} to[bend left=14] (u) to[bend right=13] (y) node[left] {$x(u,v';\alpha)$};
\draw (.75,.82) to[bend left=23] (.80,1.23) node[pos=.45,right] {$\theta$};
\draw (1.1,1.45) to[bend left=18] (1.62,1.52) node[pos=.53,above] {$\alpha$};
\draw (1.19,1.52) to[bend left=17] (1.60,1.61);
\draw (1.56,2.48) to[bend left=12] (1.60,2.92) node[pos=.55,left] {$\theta$};
\end{tikzpicture}}

\[
\resizebox{\linewidth}{!}{$\displaystyle
\left\{\begin{aligned}
\|x(u,v;\alpha)-u\|^2 &\underset{\substack{\text{par raison}\\\text{de sym.}}}{=}
\|x(u,v',\alpha)-u\|^2=(\cos\alpha-1)^2+\sin^2\alpha=2(1-\cos\alpha)=4\sin^2\frac{\alpha}{2}\\
\|x(u,v,\alpha)-x(u,v',\alpha)\|^2&=\sin^2\alpha\,\|v-v'\|^2=\sin^2\alpha\,4\sin^2\frac{\theta}{2}
=\left(4\sin^2\frac{\alpha}{2}\cos^2\frac{\alpha}{2}\right)\left(4\sin^2\frac{\theta}{2}\right)
\end{aligned}\right.$}
\]
% archive sheet number: 3
